\documentclass[12pt,a4paper]{article}

% ─── Packages ───────────────────────────────────────────────────────────────
\usepackage{amsmath,amssymb,amsthm}
\usepackage{hyperref}
\usepackage{geometry}
\usepackage{microtype}
\usepackage{booktabs}
\usepackage{array}
\usepackage{enumitem}
\usepackage{xcolor}

\geometry{margin=2.5cm}
\hypersetup{colorlinks=true, linkcolor=blue, citecolor=blue, urlcolor=blue}

% ─── Theorem environments ────────────────────────────────────────────────────
\newtheorem{theorem}{Theorem}[section]
\newtheorem{lemma}[theorem]{Lemma}
\newtheorem{corollary}[theorem]{Corollary}
\newtheorem{proposition}[theorem]{Proposition}
\theoremstyle{definition}
\newtheorem{definition}[theorem]{Definition}
\newtheorem{example}[theorem]{Example}
\theoremstyle{remark}
\newtheorem{remark}[theorem]{Remark}
\newtheorem{conjecture}[theorem]{Conjecture}

% ─── Custom commands ─────────────────────────────────────────────────────────
\newcommand{\frac}{\operatorname{frac}}   % fractional part
\newcommand{\Z}{\mathbb{Z}}
\newcommand{\N}{\mathbb{N}}
\newcommand{\R}{\mathbb{R}}
\newcommand{\Kex}{K_{\mathrm{exact}}}
\newcommand{\dK}{d(K)}

% ─────────────────────────────────────────────────────────────────────────────
\begin{document}

\title{\textbf{Wieferich Primes as Zeros of a Sawtooth Function}\\[6pt]
{\large A New Characterization via Diophantine Iteration}}

\author{V\'ictor Manzanares Alberola\\[4pt]
\small Independent researcher, Valencia, Spain\\
\small \texttt{[contact email]}}

\date{2026}

\maketitle

% ─── Abstract ────────────────────────────────────────────────────────────────
\begin{abstract}
We introduce a real-valued sawtooth function $d(K) = \{p(K)\}$, where
$p(K)$ is the unique real solution to the equation $2^p = Kp^2 + 2$ for
even integer $K$, and $\{\cdot\}$ denotes the fractional part. We prove
that the zeros of $d(K)$ over even integers coincide exactly with the
generalized Wieferich condition $p^2 \mid (2^p - 2)$, and demonstrate
that $K$ is always even at these zeros (Theorem~\ref{thm:parity}).
We further show that iterating $K$ by steps of $2$ with a four-point
velocity extrapolation constitutes a practical algorithmic detector for
Wieferich primes. Analysing $d(K)$ across multiple bases, we prove
that the two known Wieferich primes $1093$ and $3511$ are Wieferich in
exactly the same five bases $\{2,4,8,16,32\}$---all powers of $2$---so
their apparent coincidence in five bases reduces to a single arithmetic
condition. We identify $p = 13$ as the smallest prime that is Wieferich
in three mutually independent bases. The reformulation connects to the
author's prior work on the $6k{\pm}1$ factoriser and the Kinematic
Diophantine Method, suggesting that the sawtooth structure is a special
case of a more general mechanism operating on decimal velocities of
parametric Diophantine equations.
\end{abstract}

\tableofcontents
\bigskip

% ─────────────────────────────────────────────────────────────────────────────
\section{Introduction}
% ─────────────────────────────────────────────────────────────────────────────

A prime $p$ is called a \emph{Wieferich prime in base $b$} if
$p^2 \mid (b^p - b)$.
In base $2$, only two such primes are known: $1093$ (Meissner, 1913)
and $3511$ (Beeger, 1922). Whether infinitely many exist remains one of
the deepest open problems in number theory.

The standard approach studies the Fermat quotient
$q_p(2) = (2^{p-1}-1)/p \pmod{p}$, with Wieferich primes being those
where $p \mid q_p(2)$. Here we propose a complementary
\emph{geometric} reformulation: we define a continuous real function
whose zeros over even integers are precisely the Wieferich primes
(Section~\ref{sec:sawtooth}). This transforms an arithmetic question
into a question about zeros of an explicit sawtooth function.

The approach is part of a broader programme---the
\emph{Kinematic Diophantine Method} (MDC, \cite{VMA2026MDC})---in which
integer solutions of Diophantine equations are located by analysing the
velocity, acceleration and jerk of the fractional part of a real
parametric function. The Wieferich problem and the factorisation
problem are two instances of the same kinematic principle with
different targets $\delta \in \{0, 0.5\}$.

\subsection*{Main results}

\begin{enumerate}[label=(\roman*)]
  \item \emph{Theorem~\ref{thm:zeros}}: $d(K)=0$ for even integer $K$
        if and only if $K = (2^p-2)/p^2$ for a Wieferich prime $p$.
  \item \emph{Theorem~\ref{thm:parity}}: $K$ is always even at such zeros.
  \item \emph{Theorem~\ref{thm:powers}}: the Wieferich property of
        $1093$ and $3511$ in bases $\{2,4,8,16,32\}$ is a single
        condition.
  \item \emph{Algorithm~\ref{alg:4point}}: a four-point extrapolation
        predicts $K$ at each level with precision $< 10^{-4}$ in all
        tested cases.
\end{enumerate}

% ─────────────────────────────────────────────────────────────────────────────
\section{The Sawtooth Function \texorpdfstring{$d(K)$}{d(K)}}
\label{sec:sawtooth}
% ─────────────────────────────────────────────────────────────────────────────

\subsection{Definition}

\begin{definition}
For each even positive integer $K$, let $p(K)$ be the unique real
solution $p > 2$ to
\begin{equation}\label{eq:defining}
  2^p \;=\; K p^2 + 2.
\end{equation}
Such a solution exists and is unique because $2^p$ grows faster than
$Kp^2 + 2$ and the functions cross exactly once for $p > 2$. Define
\begin{equation}\label{eq:sawtooth}
  d(K) \;=\; \{p(K)\} \;=\; p(K) - \lfloor p(K)\rfloor \;\in\; [0,1).
\end{equation}
\end{definition}

\begin{remark}
As $K$ increases by $2$, $p(K)$ increases slowly. Within each
\emph{level} (the interval where $\lfloor p(K)\rfloor = n$ for a fixed
integer $n$), $d(K)$ rises from near $0$ to near $1$, then
\emph{resets} when $p(K)$ crosses $n+1$. This gives $d(K)$ a sawtooth
shape. The local velocity is approximately
\[
  \frac{\Delta d}{\Delta K} \;\approx\;
  \frac{1}{p \ln 2 \cdot p^2/2^p},
\]
which decreases as $K$ grows, making successive teeth progressively
finer.
\end{remark}

\subsection{Zeros of \texorpdfstring{$d(K)$}{d(K)}}

\begin{theorem}[Zeros are Wieferich conditions]
\label{thm:zeros}
$d(K) = 0$ for an even integer $K > 0$ if and only if
$K = (2^p - 2)/p^2$ for some prime $p$ satisfying $p^2 \mid (2^p-2)$,
i.e., $p$ is a Wieferich prime in base $2$.
\end{theorem}

\begin{proof}
$d(K)=0$ if and only if $p(K)$ is an integer $n$, which by
\eqref{eq:defining} is equivalent to
\[
  2^n = Kn^2 + 2
  \;\iff\;
  K = \frac{2^n - 2}{n^2}.
\]
For this to be a positive even integer we need $n^2 \mid (2^n-2)$ (the
Wieferich condition when $n$ is prime) and evenness (proved in
Theorem~\ref{thm:parity}). For composite $n$, no example satisfying
$n^2 \mid (2^n-2)$ has been found for $n \le 2000$; we do not claim a
proof for this case here.
\end{proof}

% ─────────────────────────────────────────────────────────────────────────────
\section{Parity of \texorpdfstring{$K$}{K} at Zeros}
\label{sec:parity}
% ─────────────────────────────────────────────────────────────────────────────

\begin{theorem}[Parity of $K$]
\label{thm:parity}
For any odd prime $p$ and any base $b$ with $p \nmid b$, if
$K = (b^p - b)/p^2 \in \Z$ then $K$ is even.
\end{theorem}

\begin{proof}
Write $K = b(b^{p-1}-1)/p^2$.
\begin{itemize}[leftmargin=2em]
  \item \textbf{Case $b$ odd.} Then $b^{p-1} = (\text{odd})^{\text{even}}$
        is odd, so $b^{p-1}-1$ is even, hence $2 \mid b(b^{p-1}-1)$.
        Since $p$ is an odd prime, $p^2$ is odd, so $2 \mid K$.
  \item \textbf{Case $b$ even.} Then $b^p - b = b(b^{p-1}-1)$ where $b$
        is even, so $2 \mid b^p - b$, giving $2 \mid K$.
\end{itemize}
In both cases $K$ is even.
\end{proof}

\begin{corollary}
Iterating $K$ by steps of $2$ (i.e., considering only even $K$) is
sufficient: every zero of $d(K)$ over positive integers lies at an
even value.
\end{corollary}

% ─────────────────────────────────────────────────────────────────────────────
\section{Four-Point Velocity Extrapolation}
\label{sec:algorithm}
% ─────────────────────────────────────────────────────────────────────────────

Within each level $n$, the function $d(K)$ rises approximately linearly
in $K$. This quasi-linearity enables a direct prediction of $\Kex(n)$
from four consecutive evaluations.

\begin{algorithm}[Four-point Wieferich detector]
\label{alg:4point}
\mbox{}
\begin{enumerate}[label=\arabic*.]
  \item Fix $K_0$ near the start of level $n$ (i.e., $K_0 \approx 2^n/n^2$).
  \item Compute $d_i = d(K_0 + 2i)$ for $i = 0,1,2,3$.
  \item Velocities: $v_i = d_{i+1} - d_i$ for $i=0,1,2$.
  \item Mean velocity: $\bar{v} = (v_0+v_1+v_2)/3$.
  \item Predicted zero: $K_{\mathrm{pred}} = K_0 - 2\,d_0/\bar{v}$.
  \item If $K_{\mathrm{pred}} \in 2\Z$: candidate Wieferich prime $p=n$.
        Verify via modular arithmetic.
\end{enumerate}
\end{algorithm}

\subsection{Numerical validation}

Table~\ref{tab:kpred} shows the predicted and exact $K$ values at
several levels. The relative error stays below $10^{-4}$ in all cases.

\begin{table}[h]
\centering
\caption{Comparison of $K_{\mathrm{pred}}$ vs $K_{\mathrm{exact}} = (2^p-2)/p^2$}
\label{tab:kpred}
\begin{tabular}{@{}rrrrc@{}}
\toprule
$p$ & $K_{\mathrm{exact}}$ & $K_{\mathrm{pred}}$ & Error & Wieferich? \\
\midrule
$17$   & $453.5294\ldots$       & $453.53$    & $<10^{-4}$ & No \\
$19$   & $1452.316\ldots$       & $1452.31$   & $<10^{-4}$ & No \\
$23$   & $15857.478\ldots$      & $15857.48$  & $<10^{-4}$ & No \\
$1093$ & $(2^{1093}-2)/1093^2$  & exact$^*$   & $0$        & \textbf{Yes} \\
$3511$ & $(2^{3511}-2)/3511^2$  & exact$^*$   & $0$        & \textbf{Yes} \\
\bottomrule
\end{tabular}
\smallskip

{\small $^*$With sufficient precision arithmetic; prediction is exact
by construction.}
\end{table}

% ─────────────────────────────────────────────────────────────────────────────
\section{Structure Across Multiple Bases}
\label{sec:multibases}
% ─────────────────────────────────────────────────────────────────────────────

We extend the framework to general base $b$, defining
$K_b(p) = (b^p - b)/p^2$ and studying integrality across
$b \in \{2,\ldots,49\}$ and primes $p < 10^4$.

\subsection{The 1093--3511 coincidence}

\begin{theorem}[Powers-of-2 bases are a single condition]
\label{thm:powers}
If $p^2 \mid (2^p - 2)$, then $p^2 \mid ((2^k)^p - 2^k)$ for every
positive integer $k$.
\end{theorem}

\begin{proof}
$(2^k)^p - 2^k = 2^k(2^{k(p-1)}-1)$.
Write $2^{k(p-1)}-1 = (2^{p-1}-1)\bigl(1 + 2^{p-1} + \cdots +
2^{(k-1)(p-1)}\bigr)$. The factor $2^{p-1}-1$ is divisible by $p^2$
(Wieferich condition in base $2$), so $p^2 \mid (2^k)^p - 2^k$.
\end{proof}

\begin{corollary}
The fact that $1093$ and $3511$ are Wieferich in bases
$\{2,4,8,16,32\} = \{2^1,2^2,2^3,2^4,2^5\}$ is a \emph{single}
arithmetic condition---their base-$2$ Wieferich property---not five
independent coincidences.
\end{corollary}

\subsection{Independent multi-base Wieferich primes}

Modulo the dependence of Theorem~\ref{thm:powers}, we identify:

\begin{itemize}
  \item $p = 2$: Wieferich in all $b \equiv 1 \pmod{4}$ up to $49$,
        since $4 \mid b^2 - b = b(b-1)$ whenever $b \equiv 1\pmod{4}$
        (note $b-1 \equiv 0 \pmod 4$).
  \item $p = 3$: Wieferich in $10$ independent bases in $[2,49]$.
  \item $p = 13$: Wieferich in bases $\{19, 22, 23\}$---three mutually
        independent bases; this is the smallest prime with $\ge 3$
        independent Wieferich bases.
  \item $p = 1093$, $p = 3511$: effectively one independent base each
        (powers of $2$).
\end{itemize}

% ─────────────────────────────────────────────────────────────────────────────
\section{Level Transition Growth Rate}
\label{sec:growth}
% ─────────────────────────────────────────────────────────────────────────────

The values $K_{\mathrm{start}}(n)$ where $p(K)$ first exceeds $n$ satisfy
$K_{\mathrm{start}}(n) \approx 2^n/n^2$. The ratio
\[
  \frac{K_{\mathrm{start}}(n+1)}{K_{\mathrm{start}}(n)}
  \;\approx\;
  2\cdot\left(\frac{n}{n+1}\right)^2 \;\longrightarrow\; 2
  \quad\text{as } n\to\infty,
\]
with convergence from below near $1.83$ for moderate $n$, consistent
with the quadratic correction.

Within level $n$, the number of even $K$ values with $d(K) < \varepsilon$
grows as $\varepsilon \cdot n^2 \cdot 2^n/2$, proportional to
$K_{\mathrm{start}}(n)$. Thus the density of near-misses per unit $K$
stays roughly constant while the spacing between levels grows
exponentially.

% ─────────────────────────────────────────────────────────────────────────────
\section{Connection to the Kinematic Diophantine Method}
\label{sec:mdc}
% ─────────────────────────────────────────────────────────────────────────────

The Wieferich detection algorithm is a special case of the
\emph{Kinematic Diophantine Method} (MDC, \cite{VMA2026MDC}), a unified
procedure for locating integer solutions of Diophantine equations
$F(x,t)=0$ by analysing the kinematics of $d(t) = \{g(x,t)\}$.

Table~\ref{tab:mdc} compares the two applications.

\begin{table}[h]
\centering
\caption{MDC unification: factorisation vs Wieferich detection}
\label{tab:mdc}
\begin{tabular}{@{}lll@{}}
\toprule
& \textbf{Factorisation} & \textbf{Wieferich} \\
\midrule
Parameter $t$      & $m$ (factor candidate)     & $K$ (even integer) \\
Function $g(t)$    & $N/(2(2m+3))$              & $p(K)$: solves $2^p=Kp^2+2$ \\
Target $\delta$    & $0.5$                      & $0.0$ \\
Search space       & $L_1 = \{6k\pm 1\}$        & $2\Z$ \\
Solution           & factor $f = 2m+3$ of $N$   & Wieferich prime $p = n$ \\
Step size          & $1$                        & $2$ \\
Cost per jump      & $O(1)$ divisions           & $O(1)$ Newton steps \\
\bottomrule
\end{tabular}
\end{table}

In both cases, $d(t)$ is a real sawtooth whose local kinematics allow
jumping directly to the solution. The step-by-$2$ constraint in the
Wieferich case and the $L_1$ restriction in the factorisation case are
consequences of the underlying modular arithmetic, not ad hoc choices.

% ─────────────────────────────────────────────────────────────────────────────
\section{Conjectures}
\label{sec:conjectures}
% ─────────────────────────────────────────────────────────────────────────────

\begin{conjecture}\label{conj:infinitely-many}
The function $d(K)$ has infinitely many zeros over even integers.
Equivalently, there exist infinitely many Wieferich primes in base $2$.
\end{conjecture}

\begin{conjecture}\label{conj:no-composite}
No composite $n$ satisfies $n^2 \mid (2^n - 2)$. Equivalently, all
zeros of $d(K)$ over even integers correspond to prime values of $p(K)$.
\end{conjecture}

\begin{conjecture}\label{conj:mod6}
For every prime $p > 3$, if $K_b(p) = (b^p - b)/p^2 \in \Z$ for some
base $b$ with $p \nmid 6b$, then $K_b(p) \equiv 0 \pmod{6}$.
\end{conjecture}

% ─────────────────────────────────────────────────────────────────────────────
\section{Conclusion}
\label{sec:conclusion}
% ─────────────────────────────────────────────────────────────────────────────

We have introduced a sawtooth function $d(K)$ whose zeros over even
integers characterise Wieferich primes. The principal contributions are:

\begin{enumerate}[label=(\roman*)]
  \item A geometric reformulation of the Wieferich condition as zeros of
        an explicit real function (Theorem~\ref{thm:zeros}).
  \item A proof that $K$ is always even at these zeros
        (Theorem~\ref{thm:parity}), explaining why stepping by $2$
        loses no information.
  \item A practical four-point velocity algorithm (Algorithm~\ref{alg:4point})
        that predicts $\Kex(n)$ with relative error $< 10^{-4}$.
  \item A structural explanation of the $1093$--$3511$ coincidence
        (Theorem~\ref{thm:powers}).
  \item Identification of $p=13$ as the smallest prime with three
        independent Wieferich bases.
\end{enumerate}

The connection to the MDC framework and the $6k\pm 1$ sieve suggests
that the sawtooth mechanism is part of a broader structure in which
decimal velocities of parametric Diophantine equations encode
arithmetic properties of integers in a continuous, geometrically
accessible form.

% ─────────────────────────────────────────────────────────────────────────────
\section*{Acknowledgements}
% ─────────────────────────────────────────────────────────────────────────────

The author thanks Claude (Anthropic) for computational assistance and
interactive discussion during the development of these results
(conversation of 26 February 2026).

% ─────────────────────────────────────────────────────────────────────────────
\begin{thebibliography}{9}

\bibitem{Meissner1913}
W.~Meissner,
\textit{\"Uber die Teilbarkeit von $2^{p-1}-1$ durch das Quadrat der
Primzahl $p=1093$},
Sitzungsber. Akad. Berlin, 1913, 663--667.

\bibitem{Beeger1922}
N.~G.~W.~H.~Beeger,
\textit{On a new case of the congruence $2^{p-1} \equiv 1\pmod{p^2}$},
Messenger Math. 51 (1922), 149--150.

\bibitem{Ribenboim1996}
P.~Ribenboim,
\textit{The New Book of Prime Number Records},
Springer, 1996.

\bibitem{McIntosh2004}
R.~J.~McIntosh and E.~L.~Roettger,
\textit{A search for Fibonacci-Wieferich and Wolstenholme primes},
Math.\ Comp.\ 76 (2007), 2087--2094.

\bibitem{VMA2026MDC}
V.~Manzanares Alberola,
\textit{M\'etodo Diof\'antico Cinem\'atico: aplicaciones a
factorizaci\'on y primos de Wieferich},
Preprint, Valencia, 2026.

\bibitem{VMA2025}
V.~Manzanares Alberola,
\textit{Prime Number Theory: Modular Methods and Density Models},
Preprint, Valencia, 2025.
\texttt{https://github.com/[username]/prime-modular-methods}

\end{thebibliography}

\end{document}
